\section{Frozen intervention-policy protocol}
\label{supp:policy_protocol}
We turn the audit into a frozen prompt-level decision rule and evaluate the
selected configurations on new samples. For each of 128 held-out prompts,
eight diagnosis seeds compare two numerical candidates with the reference.
A candidate is eligible when mean predicted preference is at least 0.475,
CLIP alignment loss is at most 0.01, DINO and CLIP spread ratios are at least
0.9, DINO coverage is at least 0.8, and both support-tail excesses are at most
0.05. Among eligible candidates, the rule chooses higher DINO coverage, then
higher preference, then a fixed solver order. It otherwise uses the frozen
higher-budget fallback. These small-sample point criteria select an action;
independent outcomes determine whether that action is adequate.

The SD1.5 shortlist is DPM++8/UniPC8 with DPM++16 fallback. The flow transfer
uses DPM++8/Euler8 with native20 fallback and exactly the same numerical
thresholds. Reference fit and calibration are model-specific. Sixteen separate
outcome seeds per prompt evaluate each policy. An approximate fixed-policy intersection--union
test\cite{berger1996intersection} requires all one-sided 95\% prompt-bootstrap constraints to pass:
preference at least 0.45, alignment loss at most 0.01, both spread ratios at
least 0.9, and both support excesses at most 0.05. The fallback receives the
same evaluation. All-arm intervals are descriptive; four prespecified paired
policy contrasts use Bonferroni intervals within each study. The original
coverage criterion uses shared initial-noise seeds. The secondary
independent-noise sensitivity specified after outcome access below evaluates the same saved decisions with
disjoint outcome halves; it does not retrospectively change eligibility or gates.

The comparators include fixed strong solvers, preference-and-alignment selection,
quality-constrained Vendi, Vendi alone, PRDC harmonic selection, a random
shortlist of the same size, and exhaustive search over the declared eight-step
candidate pool. The full matrix is generated for fair comparison, while policy
costs count only the candidates observed by that policy. A two-candidate audit
requires $8(50+8+8)=528$ denoiser calls per prompt before outcome generation,
plus feature extraction and scoring. Fixed solvers have no online diagnostic
cost. We report inference savings and their amortization separately from
this overhead.
Preparing the support audit also generates $(32\times16+256)\times50=38{,}400$
reference calls per family for its named policy reference. This reusable
fitting/calibration pool is separate from per-prompt diagnosis. Charging it
across the 128 evaluation prompts adds 300 calls per prompt: the lower-bound
cold-start break-even is then 428 outputs for SD1.5 and 145 for Sana, rather
than the prepared-audit bounds of 273 and 93. All bounds exclude feature and
reward computation. Additional reference pools, development runs, and the full
comparison matrix belong to the research cost ledger rather than a claim of
realized deployment savings.
An additional prespecified assignment control randomly permutes the audit's
selected actions across prompts, preserving their exact counts and inference
cost. It tests prompt-specific assignment using independent outcome preference
and the fraction of prompts whose observed 16-seed preference falls below 0.4.
This secondary replay receives the same diagnostic cost; its conditional
bootstrap intervals are approximate because the permutation couples prompts.
\section{Complete intervention-policy outcomes}
\label{supp:policy_outcomes}
Every frozen policy is evaluated with outcome replicas 8--23. Coverage in
the original tables compares coupled initial-noise samples; the independent-
noise sensitivity below keeps cardinality fixed and reports its different
estimand separately. The four primary contrasts retain their original within-study multiplicity adjustment. The two secondary assignment contrasts are a separate family. No policy is chosen from these outcomes.
\paragraph{SD1.5.}
The original policy reaches learned preference 0.482 at 14.06 calls per image.
The complete absolute outcomes and adjusted contrasts are reported below.
\begin{table*}[t]\centering\small
\caption{All SD1.5 policies. Spread ratios are ratios of mean within-prompt variances. Intervals for all policies are descriptive; only the frozen audit receives the primary intersection--union interpretation. Inc.: inconclusive.}
\begin{tabular}{lcccccc}\toprule
Policy & Calls & $\psi$ & DINO spread & CLIP spread & Paired cov. & Gate \\\midrule
Audit & 14.06 & 0.482 [0.477, 0.488] & 1.000 [0.988, 1.010] & 1.003 [0.987, 1.019] & 0.979 & pass \\
Preference + alignment & 12.62 & 0.477 [0.472, 0.483] & 1.008 [0.994, 1.021] & 1.012 [0.995, 1.030] & 0.970 & pass \\
Quality + Vendi & 12.62 & 0.477 [0.472, 0.482] & 1.009 [0.996, 1.022] & 1.010 [0.993, 1.029] & 0.968 & pass \\
Vendi only & 8.00 & 0.444 [0.436, 0.452] & 1.022 [1.005, 1.039] & 1.038 [1.017, 1.061] & 0.945 & inc. \\
PRDC harmonic & 8.00 & 0.447 [0.439, 0.455] & 1.017 [1.001, 1.034] & 1.029 [1.008, 1.052] & 0.944 & inc. \\
Random shortlist & 15.25 & 0.487 [0.482, 0.492] & 1.001 [0.987, 1.017] & 1.015 [0.998, 1.034] & 0.984 & pass \\
Exhaustive & 13.38 & 0.482 [0.477, 0.488] & 1.007 [0.994, 1.019] & 1.002 [0.984, 1.020] & 0.968 & pass \\
Fixed DPM++8 & 8.00 & 0.448 [0.440, 0.456] & 1.014 [0.998, 1.032] & 1.031 [1.011, 1.054] & 0.949 & inc. \\
Fixed UniPC8 & 8.00 & 0.440 [0.432, 0.448] & 1.027 [1.011, 1.045] & 1.037 [1.016, 1.061] & 0.939 & fail \\
Fixed DPM++16 & 16.00 & 0.490 [0.485, 0.494] & 0.994 [0.983, 1.005] & 1.010 [0.996, 1.025] & 0.990 & pass \\
Fixed UniPC16 & 16.00 & 0.488 [0.483, 0.493] & 1.004 [0.992, 1.016] & 1.021 [1.006, 1.036] & 0.987 & pass \\
Shuffled audit & 14.06 & 0.476 [0.469, 0.483] & 1.005 [0.992, 1.018] & 1.021 [1.005, 1.038] & 0.975 & pass \\
\bottomrule\end{tabular}\end{table*}
\begin{table}[t]\centering\small
\caption{SD1.5 audit-minus-baseline contrasts. Intervals are Bonferroni-adjusted within the four primary or two secondary comparisons. Assignment-control inference is approximate conditional on the realized permutation.}
\begin{tabular}{p{.52\columnwidth}r}\toprule
Contrast & Difference [interval] \\\midrule
Preference vs. DPM++8 & 0.034 [0.025, 0.044] \\
Calls vs. DPM++16 & -1.938 [-2.688, -1.250] \\
Calls vs. random shortlist & -1.188 [-1.938, -0.500] \\
Coverage vs. quality + Vendi & 0.011 [0.003, 0.021] \\
\midrule
Preference vs. shuffled audit & 0.007 [-0.001, 0.014] \\
Low-preference rate vs. shuffled audit & -0.039 [-0.094, 0.008] \\
\bottomrule\end{tabular}\end{table}
\paragraph{Sana.}
The unchanged decision rule reaches learned preference 0.502 at 14.28 calls
per image. Its stronger 20-call fallback remains an explicit comparator.
\begin{table*}[t]\centering\small
\caption{All Sana policies. Spread ratios are ratios of mean within-prompt variances. Intervals for all policies are descriptive; only the frozen audit receives the primary intersection--union interpretation. Inc.: inconclusive.}
\begin{tabular}{lcccccc}\toprule
Policy & Calls & $\psi$ & DINO spread & CLIP spread & Paired cov. & Gate \\\midrule
Audit & 14.28 & 0.502 [0.498, 0.506] & 1.043 [1.022, 1.066] & 1.076 [1.047, 1.107] & 0.950 & pass \\
Preference + alignment & 11.94 & 0.505 [0.500, 0.510] & 1.053 [1.028, 1.080] & 1.074 [1.040, 1.113] & 0.910 & pass \\
Quality + Vendi & 11.94 & 0.504 [0.499, 0.509] & 1.060 [1.034, 1.088] & 1.083 [1.048, 1.122] & 0.903 & pass \\
Vendi only & 8.00 & 0.478 [0.470, 0.486] & 1.140 [1.102, 1.181] & 1.158 [1.110, 1.210] & 0.812 & pass \\
PRDC harmonic & 8.00 & 0.480 [0.472, 0.488] & 1.127 [1.093, 1.166] & 1.159 [1.119, 1.204] & 0.876 & pass \\
Random shortlist & 14.84 & 0.503 [0.499, 0.507] & 1.040 [1.020, 1.063] & 1.072 [1.042, 1.106] & 0.953 & pass \\
Exhaustive & 14.28 & 0.502 [0.498, 0.506] & 1.043 [1.022, 1.066] & 1.076 [1.047, 1.107] & 0.950 & pass \\
Fixed DPM++8 & 8.00 & 0.479 [0.471, 0.487] & 1.134 [1.097, 1.174] & 1.164 [1.123, 1.209] & 0.882 & pass \\
Fixed Euler8 & 8.00 & 0.475 [0.466, 0.484] & 1.118 [1.080, 1.158] & 1.146 [1.097, 1.203] & 0.761 & pass \\
Fixed native20 & 20.00 & 0.510 [0.508, 0.513] & 1.001 [0.992, 1.010] & 1.018 [1.002, 1.034] & 0.998 & pass \\
Shuffled audit & 14.28 & 0.495 [0.488, 0.502] & 1.060 [1.031, 1.092] & 1.091 [1.058, 1.128] & 0.939 & pass \\
\bottomrule\end{tabular}\end{table*}
\begin{table}[t]\centering\small
\caption{Sana audit-minus-baseline contrasts. Intervals are Bonferroni-adjusted within the four primary or two secondary comparisons. Assignment-control inference is approximate conditional on the realized permutation.}
\begin{tabular}{p{.52\columnwidth}r}\toprule
Contrast & Difference [interval] \\\midrule
Preference vs. DPM++8 & 0.023 [0.013, 0.034] \\
Calls vs. native20 & -5.719 [-7.031, -4.312] \\
Calls vs. random shortlist & -0.562 [-1.219, -0.094] \\
Coverage vs. quality + Vendi & 0.047 [0.018, 0.081] \\
\midrule
Preference vs. shuffled audit & 0.007 [-0.002, 0.017] \\
Low-preference rate vs. shuffled audit & -0.047 [-0.102, 0.000] \\
\bottomrule\end{tabular}\end{table}

\begin{table*}[t]
\centering
\small
\setlength{\tabcolsep}{3.4pt}
\caption{Operating points of unchanged policies. Absolute outcomes
show saved descriptive 95\% paired prompt-bootstrap intervals over 128
prompts. Preference excludes shared-noise pairs; DINO coverage uses disjoint
eight-sample halves, averaged in both directions. The contrast block uses
99.5\% preference intervals (ten comparisons) and 99.375\% coverage intervals
(eight comparisons). These sensitivities were specified after access to
original outcomes. Inference calls exclude 528 diagnosis calls/prompt and
reusable preparation; fallbacks are SD DPM++16 and Sana native20.}
\label{tab:policy}
\begin{tabular*}{\textwidth}{@{\extracolsep{\fill}}llrrrrr@{}}
\toprule
& & & \multicolumn{2}{c}{Learned preference $\uparrow$}
    & \multicolumn{2}{c@{}}{DINO coverage $\uparrow$} \\
\cmidrule(lr){4-5}\cmidrule(l){6-7}
Model & Policy or comparator & Calls/image & Estimate & Interval & Estimate & Interval \\
\midrule
\multicolumn{7}{@{}l}{\textit{Absolute outcomes}} \\
SD1.5 & \textbf{Audit} & $14.06$ & $0.484$ & $[0.479,\,0.489]$ & $0.913$ & $[0.897,\,0.928]$ \\
 & Quality + Vendi & $12.62$ & $0.479$ & $[0.474,\,0.484]$ & $0.902$ & $[0.883,\,0.920]$ \\
 & Same-count shuffle & $14.06$ & $0.478$ & $[0.471,\,0.484]$ & $0.912$ & $[0.897,\,0.927]$ \\
 & Fixed DPM++8 & $8.00$ & $0.451$ & $[0.444,\,0.459]$ & $0.895$ & $[0.876,\,0.913]$ \\
 & Fallback (DPM++16) & $16.00$ & $0.491$ & $[0.486,\,0.495]$ & $0.927$ & $[0.914,\,0.940]$ \\
\addlinespace[2pt]
Sana & \textbf{Audit} & $14.28$ & $0.502$ & $[0.498,\,0.506]$ & $0.840$ & $[0.810,\,0.867]$ \\
 & Quality + Vendi & $11.94$ & $0.503$ & $[0.498,\,0.508]$ & $0.789$ & $[0.750,\,0.825]$ \\
 & Same-count shuffle & $14.28$ & $0.495$ & $[0.488,\,0.502]$ & $0.844$ & $[0.816,\,0.870]$ \\
 & Fixed DPM++8 & $8.00$ & $0.480$ & $[0.472,\,0.487]$ & $0.775$ & $[0.743,\,0.807]$ \\
 & Fallback (native20) & $20.00$ & $0.510$ & $[0.507,\,0.512]$ & $0.906$ & $[0.889,\,0.922]$ \\
\midrule
\multicolumn{7}{@{}l}{\textit{Audit minus comparator}} \\
SD1.5 & Quality + Vendi & $+1.44$ & $+0.0053$ & $[0.0019,\,0.0097]$ & $+0.0107$ & $[-0.0015,\,0.0269]$ \\
 & Same-count shuffle & $0.00$ & $+0.0063$ & $[-0.0025,\,0.0150]$ & $+0.0010$ & $[-0.0132,\,0.0161]$ \\
Sana & Quality + Vendi & $+2.34$ & $-0.0014$ & $[-0.0057,\,0.0024]$ & $+0.0508$ & $[0.0148,\,0.0947]$ \\
 & Same-count shuffle & $0.00$ & $+0.0068$ & $[-0.0038,\,0.0190]$ & $-0.0039$ & $[-0.0453,\,0.0365]$ \\
\bottomrule
\end{tabular*}
\end{table*}


\subsection{Sensitivity to sampling coupling}
\label{sec:t2i-coupling-sensitivity}
Two secondary sensitivity protocols were specified after access to the original outcome analyses; the original paired-seed results remain the primary record. Neither changes an intervention, routing threshold, candidate pool, or saved decision. Their uncertainty remains approximate prompt-bootstrap inference.

\paragraph{Coverage at fixed cardinality.}
For each of the 128 held-out prompts, divide the 16 outcome seeds into
$A=\{8,\ldots,15\}$ and $B=\{16,\ldots,23\}$.
Write $R_A,R_B$ for the reference samples and $Q_A,Q_B$ for candidate
samples from the two halves. Paired coverage averages $(R_A,Q_A)$ and
$(R_B,Q_B)$, while disjoint coverage averages $(R_A,Q_B)$ and $(R_B,Q_A)$.
Both use eight reference and eight candidate samples and the third-nearest
reference neighbor, preserving cardinality while changing coupling.
We retain all 29 SD1.5 and seven Sana arms, both feature representations,
all four PRDC measurements, and all 23 frozen policies. The bootstrap
resamples whole prompts, preserving dependence between the two directions.
The eight declared DINO-coverage contrasts use 99.375\% percentile
intervals with Bonferroni allocation across both families. The saved
percentile levels are .003125/.996875.

Let $\widehat C$ be the fraction of $n$ reference points whose
$k$-nearest-neighbor ball contains at least one of the $m$ candidate points,
where $1\leq k<n$. For independent reference and candidate draws from the
same continuous distribution with no distance ties, its expectation is
known~\cite[Lemma~2]{naeem2020density}:
\begin{equation}
 \mathbb E[\widehat C]=1-\prod_{\ell=1}^{k}\frac{n-\ell}{m+n-\ell}.
\label{eq:supp_expected_coverage}
\end{equation}
At $n=m=8,k=3$ it is $12/13\approx.9231$, whereas using the identical image set as both inputs gives one. Independent reference-self coverage is .9302 on SD1.5 (descriptive 95\% interval $[.9155,.9443]$) and .9097 on Sana ($[.8911,.9272]$). Thus neither a shortfall from one nor agreement with this finite-sample expectation quantifies missing population support. The observed audit--shuffle disjoint-coverage differences are .0010 on SD1.5 ($[-.0132,.0161]$) and $-.0039$ on Sana ($[-.0453,.0365]$), with the declared adjusted intervals. The SD1.5 audit--quality+Vendi interval includes zero; Sana retains a positive difference at increased cost.

\paragraph{Independent-noise preference.}
Let $a_i,b_i$ be candidate and reference PickScore logits for outcome seed $i$ at one prompt, and let $\sigma$ denote the sigmoid. The original score is $\widehat\psi_{\rm pair}=m^{-1}\sum_i\sigma(a_i-b_i)$, with $m=16$. We additionally compute
\begin{equation}
 \widehat\psi_{\rm ind}=\frac{1}{m(m-1)}\sum_{i\ne j}\sigma(a_i-b_j).
\label{eq:supp_independent_preference}
\end{equation}
This is a symmetrized order-two U-statistic over independent paired seed units and estimates preference between independent draws from the two conditional distributions. Its 240 within-prompt cross-pairs are dependent; only the 128 prompt units are bootstrapped (10,000 repetitions, fixed seed 74117). For the same reference, both estimators equal one half. Original diagonal policy means reproduce exactly. All 94 arm/reference comparisons (SD 29 against three references, Sana 7 against its reference), all 23 frozen policies, and every audit-minus-other-policy contrast are reported in the saved CSV/JSON records. No outcome is used to reroute a prompt.

The ten declared preference contrasts use 99.5\% intervals across both families. Audit independent-minus-paired preference is .00156 on SD1.5 ($[.00075,.00251]$) and $-.00020$ on Sana ($[-.00055,.00020]$). Independent-noise audit--shuffle differences remain unresolved: .00625 ($[-.00251,.01497]$) and .00685 ($[-.00384,.01900]$). Mean raw ImageReward differences are unchanged by pairing, by linearity; no rescoring is needed. These checks identify which conclusions depend on coupling without presenting the exploratory sensitivities as untouched confirmation.



\begin{table*}[tp]
\centering\small
\caption{Complete independent-noise preference sensitivity for all 23 frozen policies. The final-column intervals are descriptive 95\% paired-prompt bootstrap intervals for independent minus paired preference.}
\begin{tabular}{llrrrr}
\toprule Family & Policy & NFE & Paired & Independent & Difference [95\% CI] \\ \midrule
SD1.5 & Audit & 14.06 & 0.4825 & 0.4840 & +0.0016 [0.0010, 0.0022]\\
SD1.5 & Preference + alignment & 12.62 & 0.4774 & 0.4791 & +0.0017 [0.0011, 0.0024]\\
SD1.5 & Quality + Vendi & 12.62 & 0.4769 & 0.4787 & +0.0017 [0.0012, 0.0024]\\
SD1.5 & Vendi only & 8.00 & 0.4441 & 0.4473 & +0.0032 [0.0024, 0.0040]\\
SD1.5 & PRDC harmonic & 8.00 & 0.4469 & 0.4501 & +0.0032 [0.0024, 0.0040]\\
SD1.5 & Random shortlist & 15.25 & 0.4866 & 0.4877 & +0.0011 [0.0006, 0.0017]\\
SD1.5 & Exhaustive & 13.38 & 0.4825 & 0.4840 & +0.0016 [0.0010, 0.0022]\\
SD1.5 & Fixed DPM++8 & 8.00 & 0.4481 & 0.4513 & +0.0032 [0.0024, 0.0040]\\
SD1.5 & Fixed UniPC8 & 8.00 & 0.4403 & 0.4438 & +0.0035 [0.0027, 0.0043]\\
SD1.5 & Fixed DPM++16 & 16.00 & 0.4895 & 0.4906 & +0.0011 [0.0005, 0.0017]\\
SD1.5 & Fixed UniPC16 & 16.00 & 0.4880 & 0.4892 & +0.0011 [0.0005, 0.0018]\\
SD1.5 & Same-count shuffle & 14.06 & 0.4758 & 0.4778 & +0.0020 [0.0013, 0.0027]\\
Sana & Audit & 14.28 & 0.5023 & 0.5021 & -0.0002 [-0.0004, 0.0001]\\
Sana & Preference + alignment & 11.94 & 0.5051 & 0.5047 & -0.0004 [-0.0008, -0.0001]\\
Sana & Quality + Vendi & 11.94 & 0.5039 & 0.5035 & -0.0004 [-0.0007, -0.0001]\\
Sana & Vendi only & 8.00 & 0.4781 & 0.4782 & +0.0001 [-0.0003, 0.0005]\\
Sana & PRDC harmonic & 8.00 & 0.4800 & 0.4802 & +0.0003 [-0.0001, 0.0007]\\
Sana & Random shortlist & 14.84 & 0.5028 & 0.5026 & -0.0002 [-0.0005, 0.0001]\\
Sana & Exhaustive & 14.28 & 0.5023 & 0.5021 & -0.0002 [-0.0004, 0.0001]\\
Sana & Fixed DPM++8 & 8.00 & 0.4794 & 0.4797 & +0.0003 [-0.0001, 0.0007]\\
Sana & Fixed Euler8 & 8.00 & 0.4749 & 0.4750 & +0.0001 [-0.0004, 0.0006]\\
Sana & Fixed native20 & 20.00 & 0.5101 & 0.5097 & -0.0005 [-0.0007, -0.0003]\\
Sana & Same-count shuffle & 14.28 & 0.4954 & 0.4953 & -0.0002 [-0.0005, 0.0002]\\
\bottomrule
\end{tabular}
\end{table*}

\begin{table*}[tp]
\centering\small
\caption{All ten declared independent-noise preference contrasts. Bonferroni percentile intervals are 99.5\% each across the ten-contrast family; all policies remain unchanged.}
\begin{tabular}{llrr}
\toprule Family & Contrast & Difference & Adjusted interval \\ \midrule
SD1.5 & Audit: independent minus paired & +0.0016 & [+0.00075, +0.00251]\\
SD1.5 & Quality+Vendi: independent minus paired & +0.0017 & [+0.00094, +0.00271]\\
SD1.5 & Independent audit minus DPM++8 & +0.0327 & [+0.02320, +0.04206]\\
SD1.5 & Independent audit minus quality+Vendi & +0.0053 & [+0.00185, +0.00975]\\
SD1.5 & Independent audit minus same-count shuffle & +0.0063 & [-0.00251, +0.01497]\\
Sana & Audit: independent minus paired & -0.0002 & [-0.00055, +0.00020]\\
Sana & Quality+Vendi: independent minus paired & -0.0004 & [-0.00089, +0.00004]\\
Sana & Independent audit minus DPM++8 & +0.0225 & [+0.01167, +0.03440]\\
Sana & Independent audit minus quality+Vendi & -0.0014 & [-0.00568, +0.00235]\\
Sana & Independent audit minus same-count shuffle & +0.0068 & [-0.00384, +0.01900]\\
\bottomrule
\end{tabular}
\end{table*}

\begin{table*}[tp]
\centering\small
\caption{All frozen policies under the fixed-cardinality coverage sensitivity. DINO coverage uses eight samples in each set and $k=3$. Shared-seed and disjoint-seed comparisons each average their two outcome halves. Brackets are descriptive 95\% paired-prompt bootstrap intervals.}
\begin{tabular}{llcc}
\toprule Family & Policy & Shared-seed coverage & Disjoint-seed coverage \\ \midrule
SD1.5 & Audit & 0.991 [0.986, 0.995] & 0.913 [0.897, 0.928]\\
SD1.5 & Preference + alignment & 0.992 [0.988, 0.996] & 0.903 [0.885, 0.921]\\
SD1.5 & Quality + Vendi & 0.991 [0.986, 0.995] & 0.902 [0.883, 0.920]\\
SD1.5 & Vendi only & 0.975 [0.966, 0.982] & 0.882 [0.858, 0.903]\\
SD1.5 & PRDC harmonic & 0.974 [0.966, 0.981] & 0.888 [0.869, 0.906]\\
SD1.5 & Random shortlist & 0.993 [0.986, 0.998] & 0.921 [0.906, 0.936]\\
SD1.5 & Exhaustive & 0.985 [0.978, 0.991] & 0.908 [0.891, 0.925]\\
SD1.5 & Fixed DPM++8 & 0.983 [0.977, 0.989] & 0.895 [0.876, 0.913]\\
SD1.5 & Fixed UniPC8 & 0.970 [0.961, 0.978] & 0.866 [0.841, 0.889]\\
SD1.5 & Fixed DPM++16 & 0.997 [0.995, 0.999] & 0.927 [0.914, 0.940]\\
SD1.5 & Fixed UniPC16 & 0.996 [0.993, 0.999] & 0.914 [0.897, 0.930]\\
SD1.5 & Same-count shuffle & 0.991 [0.986, 0.995] & 0.912 [0.897, 0.927]\\
Sana & Audit & 0.968 [0.955, 0.979] & 0.840 [0.810, 0.867]\\
Sana & Preference + alignment & 0.938 [0.915, 0.958] & 0.787 [0.750, 0.820]\\
Sana & Quality + Vendi & 0.930 [0.902, 0.954] & 0.789 [0.750, 0.825]\\
Sana & Vendi only & 0.870 [0.836, 0.901] & 0.708 [0.667, 0.746]\\
Sana & PRDC harmonic & 0.928 [0.906, 0.946] & 0.772 [0.738, 0.805]\\
Sana & Random shortlist & 0.971 [0.958, 0.982] & 0.837 [0.806, 0.866]\\
Sana & Exhaustive & 0.968 [0.955, 0.979] & 0.840 [0.810, 0.867]\\
Sana & Fixed DPM++8 & 0.934 [0.917, 0.950] & 0.775 [0.743, 0.807]\\
Sana & Fixed Euler8 & 0.837 [0.800, 0.871] & 0.666 [0.624, 0.705]\\
Sana & Fixed native20 & 0.999 [0.998, 1.000] & 0.906 [0.889, 0.922]\\
Sana & Same-count shuffle & 0.964 [0.951, 0.977] & 0.844 [0.816, 0.870]\\
\bottomrule\end{tabular}
\end{table*}
\begin{table*}[tp]
\centering\small
\caption{All eight declared coverage sensitivity contrasts. Percentile intervals are 99.375\% each across the eight-comparison family. These secondary sensitivities were specified after access to earlier outcomes and evaluate unchanged policy decisions.}
\begin{tabular}{llc}
\toprule Family & Contrast & Difference [adjusted interval] \\ \midrule
SD1.5 & Audit: disjoint minus shared seeds & -0.0776 [-0.0991, -0.0571]\\
SD1.5 & Quality + Vendi: disjoint minus shared seeds & -0.0884 [-0.1143, -0.0649]\\
SD1.5 & Disjoint audit minus quality + Vendi & +0.0107 [-0.0015, +0.0269]\\
SD1.5 & Disjoint audit minus same-count shuffle & +0.0010 [-0.0132, +0.0161]\\
Sana & Audit: disjoint minus shared seeds & -0.1279 [-0.1610, -0.0973]\\
Sana & Quality + Vendi: disjoint minus shared seeds & -0.1406 [-0.1742, -0.1089]\\
Sana & Disjoint audit minus quality + Vendi & +0.0508 [+0.0148, +0.0947]\\
Sana & Disjoint audit minus same-count shuffle & -0.0039 [-0.0453, +0.0365]\\
\bottomrule\end{tabular}
\end{table*}

